\documentclass[11pt,a4paper]{article}
\usepackage[margin=1in]{geometry}
\usepackage{amsmath,amssymb,amsfonts}
\usepackage{mathtools}
\usepackage{lmodern}
\usepackage[T1]{fontenc}
\usepackage[utf8]{inputenc}
\usepackage{textcomp}
\usepackage{microtype}
\usepackage{parskip}
\usepackage{enumitem}
\usepackage{array}
\usepackage{booktabs}
\usepackage{longtable}
\usepackage{hyperref}
\hypersetup{hidelinks,
  pdftitle={Paper 047 -- The Hawking Spectrum at T_H = kappa/(2 pi) via V5 KMS},
  pdfauthor={Allen Proxmire}}
\setlength{\parindent}{0pt}
\setlength{\parskip}{6pt plus 2pt minus 1pt}

\title{\vspace{-1cm}\Large\bfseries Paper 047 --- The Hawking Spectrum at $T_H = \kappa/(2\pi)$ via V5 KMS}
\author{Allen Proxmire}
\date{May 2026}

\begin{document}
\maketitle

\section*{Preamble: What This Paper Does NOT Claim}

\begin{enumerate}[noitemsep]
\item \textbf{The 13 substrate primitives are not derived} (Paper\_087).
\item \textbf{The surface gravity $\kappa$ is not derived in this paper.} $\kappa$ is \textbf{inherited} from the GR / BH-geometric setup as the coarse-grained identification at the substrate decoupling surface (Paper\_039 \S3.2). Independent substrate-level derivation of $\kappa$ awaits the Arc ED-10 curvature-emergence program.
\item \textbf{The Euclidean conical-defect structure (\S3) is inherited via DCGT, not substrate-derived.} \S3's content (Euclidean continuation, conical-defect smoothness, imaginary-time periodicity $\beta_H = 2\pi/\kappa$) is the standard semiclassical Euclidean BH structure, applied to substrate-graph content via DCGT (Paper\_073). The substrate-level novelty is in \S6 (the cutoff corrections), not \S3.
\item \textbf{The BH mass $M$ and spacetime geometry are not derived.} Inherited from BH physics; substrate framework operates conditional on $\kappa(M, J, Q, \ldots)$.
\item \textbf{No claim that the standard semiclassical Hawking derivation is wrong.} ED's substrate-level mechanism reproduces standard Hawking at leading order ($\omega \ll \omega_c$), with \textbf{specific predicted non-thermal corrections} at higher order (under P-V5-Even).
\item \textbf{No claim that V5 supplies the only substrate-level mechanism for Hawking radiation.} Alternative substrate-level mechanisms (V1-only, or future-cascade-kernel-based) are conceivable but not load-bearing in the current substrate framework.
\item \textbf{No claim of unitarity proof.} Paper\_039 \S7 argues the information paradox is structurally not generated; this paper does not re-derive that argument. Hawking spectrum derivation operates conditional on the substrate-level structure Paper\_039 establishes.
\item \textbf{No claim that $\omega_c = c/\ell_P$ is derived.} $\omega_c$ is \textbf{chosen to match the substrate scale}: $\tau_{V5}^\mathrm{BH} = \ell_P/c$ via $\ell_\mathrm{ED} = \ell_P$ (Paper\_027 \S5.3). This is empirical-scale matching, not substrate-derivation.
\item \textbf{The $(\omega/\omega_c)^2$ leading-correction form is conditional on the P-V5-Even postulate (\S2).} For parity-asymmetric V5 envelopes, the leading correction would be $\mathcal{O}(\omega/\omega_c)$ --- linear rather than quadratic. P-V5-Even is a substrate-level structural commitment, not a derivation; alternative parity choices would predict different empirical signatures.
\item \textbf{No claim that the coefficient $c_1$ is exactly derivable.} $c_1$ is order-unity (e.g., $c_1 = 1/2$ for Gaussian V5 envelope under P-V5-Even), inherited from V5 envelope shape, which is value-layer (Paper\_090 \S3.3).
\item \textbf{Information-bearing claim is a soft structural availability, not a derivation.} \S6.3 establishes that the substrate framework leaves room for $(\omega/\omega_c)^2$ deviations to carry information; the explicit substrate-to-spectrum encoding map awaits Paper\_050.
\end{enumerate}

\section*{Abstract}

This paper delivers the \textbf{Hawking spectrum derivation} promised in Paper\_039 \S4 + \S7.2. Given the 13 postulated primitives (Paper\_087) + V1 inheritance (Paper\_089) + V5 dependence (Paper\_090) + Paper\_039's substrate-level decoupling-surface mechanism, the substrate-level Hawking spectrum follows from \textbf{V5 imaginary-time periodicity} + the \textbf{substrate-level KMS condition} at the near-horizon Euclidean geometry. The derivation has four load-bearing components: (1) V5 acquires imaginary-time periodicity $\beta_H = 2\pi/\kappa$ from the Euclidean conical-defect smoothness at the decoupling surface; (2) periodicity implies the \textbf{KMS condition} for V5 cross-chain correlators; (3) KMS implies thermal Planck distribution at temperature $T_H = 1/\beta_H = \kappa/(2\pi)$; (4) V5's substrate-level cutoff $\omega_c = c/\ell_P$ produces \textbf{specific predicted non-thermal corrections} of order $(\omega/\omega_c)^2$ (under the P-V5-Even postulate, \S2), with sign-definite and mass-dependent structure. The corrections leave \textbf{structural room} for information-bearing content, with the explicit encoding map awaiting Paper\_050 (Page-curve substrate mechanism). The headline result:
\[
n_H^\mathrm{ED}(\omega) = n_H^\mathrm{Planck}(\omega)\left[1 - c_1(\omega/\omega_c)^2 + \mathcal{O}((\omega/\omega_c)^4)\right]
\]
with $c_1$ an order-unity coefficient inherited from V5 envelope shape. The $(\omega/\omega_c)^2$ leading-correction form is \textbf{conditional on the P-V5-Even postulate (\S2)} --- V5 envelope parity-symmetric in $\omega$; parity-asymmetric envelopes would give a linear leading correction $\mathcal{O}(\omega/\omega_c)$ instead. The paper makes no claim that $\kappa$ is derived (inherited from BH geometry), no claim that $\omega_c$ is derived (chosen to match $\ell_P$), no claim that the Euclidean conical-defect content of \S3 is substrate-derived (inherited via DCGT from standard semiclassical Euclidean BH geometry), no claim of unitarity proof in this paper, and no claim that $c_1$ is exactly derivable.

\section{Introduction}

\subsection{What this paper does}

This paper supplies the \textbf{substrate-level Hawking-spectrum derivation} in four steps:

\begin{enumerate}[noitemsep]
\item V5 imaginary-time periodicity at the substrate decoupling surface (\S3).
\item KMS condition from periodicity (\S4).
\item Thermal Planck distribution at $T_H = \kappa/(2\pi)$ (\S5).
\item Non-thermal $(\omega/\omega_c)^2$ corrections from V5 finite cutoff (\S6).
\end{enumerate}

The derivation operates conditional on Paper\_039's substrate-level decoupling-surface mechanism (which establishes the substrate-graph structure at the BH horizon). This paper does not re-derive the decoupling surface; it builds the Hawking spectrum on top of it.

\subsection{Why this matters}

Paper\_039 promised the spectrum derivation in two places:

\begin{itemize}[noitemsep]
\item \S4: ``the full derivation --- including the matching to the leading-order semiclassical Hawking spectrum --- is the subject of Paper\_047.''
\item \S7.2: the non-thermal correction $n_H^\mathrm{ED} = n_H^\mathrm{Planck}[1 - c_1(\omega/\omega_c)^2 + \ldots]$ was stated as a headline result; the explicit derivation lives in this paper.
\end{itemize}

Without Paper\_047, Paper\_039's headline non-thermal-correction prediction (\S7.2) lacks dedicated derivation. This paper supplies it.

\subsection{How this fits into the BH/Hawking arc}

\begin{itemize}[noitemsep]
\item \textbf{Paper\_039:} Horizon as decoupling surface; trans-Planckian resolution; Planck-mass remnant.
\item \textbf{Paper\_047 (this paper):} Hawking spectrum via V5 KMS.
\item \textbf{Paper\_050 (in queue):} Page-curve substrate mechanism.
\item \textbf{Paper\_052 (in queue):} BH info-paradox publication-grade synthesis.
\end{itemize}

This is the second BH-arc paper; together with Paper\_039 it establishes the substrate-level Hawking mechanism.

\section{Primitive Inputs}

This paper takes the following ED substrate primitives as \textbf{postulated} (Paper\_087):

\textbf{P02, P04, P05, P07, P09, P10, P11, P12} --- as load-bearing in Paper\_039.

\textbf{V1 inheritance (Paper\_089):} finite-width retarded structure.

\textbf{V5 dependence (Paper\_090):} cross-chain correlation kernel; finite memory at near-horizon $\tau_{V5}^\mathrm{BH} \sim \ell_P/c$.

\textbf{Postulate specific to this paper:}

\textbf{P-V5-Even (V5 envelope parity-symmetric in frequency):} \emph{The V5 envelope $F_{V5}(\omega/\omega_c)$ is parity-symmetric (even) under $\omega \to -\omega$ around the substrate-saturated regime.} This postulate is required to exclude odd-power leading corrections in the $(\omega/\omega_c)$ expansion (\S6.2), giving the $(\omega/\omega_c)^2$ leading non-thermal correction. For non-even envelopes (e.g., pure exponential $e^{-\omega/\omega_c}$, which is not parity-symmetric), the leading correction is $\mathcal{O}(\omega/\omega_c)$ instead --- a different empirical signature. P-V5-Even is the substrate-level structural commitment that selects the parity-even case (Gaussian-like, $1 - (\omega/\omega_c)^2 + \ldots$); it is \textbf{not} derivable from Paper\_090's V5 admissible class alone and is labeled here for transparency.

\textbf{Upstream paper dependencies:}

\begin{itemize}[noitemsep]
\item \textbf{Paper\_087:} position paper.
\item \textbf{Paper\_089:} V1 canonical reference.
\item \textbf{Paper\_090:} V5 cross-chain kernel.
\item \textbf{Paper\_039:} Horizon decoupling --- supplies substrate-graph structure at BH horizon, including the surface gravity $\kappa$ as the coarse-grained P12-gradient identification at the decoupling surface.
\item \textbf{Paper\_073:} DCGT --- used to bridge substrate-level KMS to semiclassical Hawking continuum spectrum.
\end{itemize}

\textbf{Empirical / value-layer inputs:}

\begin{itemize}[noitemsep]
\item \textbf{Surface gravity $\kappa$:} inherited from BH geometry (Schwarzschild: $\kappa = c^4/(4GM)$; Kerr: depends on $M, J$). \textbf{Not derived in this paper.}
\item \textbf{BH mass $M$, angular momentum $J$, charge $Q$:} astrophysical parameters; inherited.
\item \textbf{Substrate scale $\ell_\mathrm{ED} = \ell_P$:} empirically identified via Paper\_027 \S5.3.
\item \textbf{V5 envelope shape:} value-layer; determines $c_1$ in non-thermal correction.
\end{itemize}

\textbf{No primitive forcing is invoked.}

\section*{\S2.5 Load-Bearing Step Audit}

{\footnotesize
\begin{longtable}{p{0.38\textwidth}p{0.16\textwidth}p{0.40\textwidth}}
\toprule
\textbf{Step} & \textbf{Status} & \textbf{Source / justification} \\
\midrule
\endhead
13 primitives postulated & P & Paper\_087 \\
V1 retarded kernel structure & D & Paper\_089 \\
V5 cross-chain correlation & D & Paper\_090 \\
Substrate decoupling surface at BH horizon & D & Paper\_039 \S3 \\
Surface gravity $\kappa$ & I & Paper\_039 \S3.2 coarse-grained from GR; this paper inherits \\
Euclidean continuation of substrate-graph structure & I-via-DCGT & \S3.1 inherited from standard semiclassical Euclidean BH geometry via DCGT (Paper\_073); substrate-level novelty is in \S6, not \S3 \\
Conical-defect smoothness condition & I-via-DCGT & \S3.2 standard Euclidean BH structure; substrate framework reinterprets the periodicity locus (V5 correlators) but inherits the geometric condition \\
Imaginary-time periodicity $\beta_H = 2\pi/\kappa$ & I-via-DCGT & \S3.3 from inherited conical-defect smoothness + inherited $\kappa$ \\
KMS condition on V5 correlators & D & \S4 from imaginary-time periodicity (substrate-level reinterpretation: periodicity now lives in V5 correlators, not in matter fields) \\
Thermal Planck distribution at $T_H = \kappa/(2\pi)$ & D & \S5 from KMS condition \\
V5 substrate cutoff $\omega_c = c/\ell_P$ & I / regime & $\tau_{V5}^\mathrm{BH} = \ell_P/c$ chosen to match substrate scale; Paper\_090 \S5.2 \\
V5 envelope parity-symmetric in $\omega$ & P (P-V5-Even) & \S2 postulate. Excludes odd-power leading corrections; required for the $(\omega/\omega_c)^2$ leading-form prediction \\
Non-thermal correction form $(\omega/\omega_c)^2$ leading & D & \S6.2 from P-V5-Even + V5 envelope Taylor expansion (even powers only) \\
Sign-definite negative leading correction & D & \S6.3 from V5 envelope decay shape (positive decay rate at large $\omega$ $\to$ negative correction to occupation number) \\
Mass-dependent scaling of correction magnitude & D & \S6.3 from $\omega_H/\omega_c \sim M_P/M$ \\
Information-bearing structure of corrections & A$\to$P & \S6.3 leaves substrate-level room for information-bearing content; explicit substrate-to-spectrum encoding map awaits Paper\_050. Soft claim, not derived \\
Coefficient $c_1$ order-unity & I & V5 envelope shape value-layer (Paper\_090 \S3.3) \\
DCGT bridge to semiclassical continuum & D & Paper\_073 \\
\bottomrule
\end{longtable}
}

All load-bearing steps are P, D, I, or labeled regime/soft. \S3 rows are labeled I-via-DCGT because they inherit from standard semiclassical Euclidean BH geometry via DCGT, not from substrate-level derivation. Even-envelope postulate (P-V5-Even) made explicit. Information-bearing claim (\S6.3) softened to A$\to$P pending Paper\_050.

\section{V5 Imaginary-Time Periodicity at the Substrate Decoupling Surface}

\subsection{Euclidean continuation of substrate-graph structure}

At the substrate decoupling surface (Paper\_039 \S3), V5 cross-chain correlations operate in the near-horizon region. To analyze thermal-like substrate behavior, we perform the \textbf{Euclidean continuation} $t \to -i\tau_E$ on the substrate-graph structure.

In standard semiclassical Hawking analysis, the Euclidean continuation of Schwarzschild spacetime near the horizon takes the form $ds^2_E = \kappa^2 r^2\,d\tau_E^2 + dr^2 + (\mathrm{angular})$ (after Rindler-like coordinates centered at the horizon). The Euclidean geometry near $r = 0$ has a \textbf{conical structure} with deficit angle determined by $\beta = 2\pi/\kappa$.

In ED's substrate framework, the substrate-graph adjacency near the decoupling surface inherits this Euclidean conical structure as the coarse-grained DCGT-continuum limit of the substrate-graph metric (Paper\_039 \S3.2). At substrate-graph scales, the structure is discrete; at coarse-grained scales, it matches the standard Euclidean BH geometry.

\subsection{Conical-defect smoothness condition}

For substrate-graph kernels to extend smoothly across the conical structure in the Euclidean continuation, the imaginary-time direction must be \textbf{periodic} with period matching the conical structure. This is the substrate-level analog of the standard semiclassical condition that the Euclidean BH geometry is regular at the horizon (no conical defect).

\textbf{The smoothness condition:} for V5 cross-chain correlators to be regular at the Euclidean horizon (no singularity in V5 envelope at the conical-defect point), the imaginary-time direction must be identified with period:
\[
\beta_H = \frac{2\pi}{\kappa}.
\]

\textbf{Dimensional check:} $\kappa$ has units $[1/\mathrm{time}]$ (acceleration / velocity = $1/\mathrm{time}$ --- surface gravity is acceleration, divided by $c$ in natural units, gives $1/\mathrm{time}$); $\beta_H$ has units $[\mathrm{time}]$. $2\pi$ dimensionless. $\checkmark$

\subsection{Imaginary-time periodicity of V5}

Combining \S3.1 and \S3.2, V5 cross-chain correlators near the decoupling surface satisfy:
\[
K_{V5}(u_A, t_A; u_B, t_B)\Big|_\mathrm{Euclidean} = K_{V5}(u_A, t_A + i\beta_H; u_B, t_B)\Big|_\mathrm{Euclidean}
\]
with periodicity $\beta_H = 2\pi/\kappa$.

This is the substrate-level imaginary-time periodicity. \textbf{It is derived from the conical-defect smoothness condition + inherited $\kappa$}, not assumed.

\section{KMS Condition from Periodicity}

\subsection{The KMS condition}

The \textbf{Kubo--Martin--Schwinger (KMS) condition} for a quantum-mechanical thermal state at temperature $T = 1/\beta$ is:
\[
\langle\phi(t_A)\phi(t_B)\rangle_T = \langle\phi(t_A + i\beta)\phi(t_B)\rangle_T
\]
for two-point correlators $\langle\phi(t_A)\phi(t_B)\rangle$ evaluated in the thermal density matrix $\rho = e^{-\beta H}/Z$.

KMS is the defining property of thermal equilibrium in quantum statistical mechanics (Kubo 1957, Martin--Schwinger 1959).

\subsection{V5 correlators satisfy KMS at $\beta = \beta_H$}

From \S3.3, V5 cross-chain correlators in the near-horizon Euclidean substrate-graph structure are periodic in imaginary time with period $\beta_H = 2\pi/\kappa$:
\[
\langle\phi(u_A, t_A)\phi(u_B, t_B)\rangle_{V5}^\mathrm{Eucl} = \langle\phi(u_A, t_A + i\beta_H)\phi(u_B, t_B)\rangle_{V5}^\mathrm{Eucl}.
\]
\textbf{This is the KMS condition} for V5 correlators at substrate level, with $\beta = \beta_H$. The substrate-level structure of V5 near the decoupling surface is \textbf{automatically thermal} at temperature
\[
T_H = \frac{1}{\beta_H} = \frac{\kappa}{2\pi}.
\]
\textbf{Dimensional check:} $T_H$ has units of $[1/\mathrm{time}]$ (in natural units with $\hbar = k_B = 1$); restoring units: $T_H = \hbar\kappa/(2\pi k_B)$ has units of temperature. $\checkmark$

\subsection{Substrate-level KMS vs.\ field-theoretic KMS}

In standard semiclassical Hawking analysis, KMS is established for \emph{matter-field} correlators on a fixed BH background --- the matter field's two-point function in the Hartle--Hawking state is KMS-thermal at $T_H$.

In ED's substrate framework, KMS is established for \textbf{V5 substrate-level cross-chain correlators} --- the substrate kernel rule-type is itself thermal at the substrate level. The two are connected via DCGT (Paper\_073 \S6.4): V5 substrate KMS coarse-grains to standard matter-field KMS in the continuum limit.

The substrate-level KMS is structurally deeper: it operates at the level of substrate primitives (V5 cross-chain kernel) rather than at the level of derived field-theoretic content.

\section{Leading-Order Hawking Spectrum}

\subsection{Thermal Planck spectrum from KMS}

The KMS condition at temperature $T_H = \kappa/(2\pi)$ (\S4.2) implies that V5 correlators in the near-horizon Euclidean substrate-graph structure have the thermal-Planck form:
\[
\langle\phi(\omega)\phi(\omega')\rangle_{V5}^\mathrm{thermal} \propto \frac{1}{e^{\hbar\omega/k_B T_H} - 1}\,\delta(\omega - \omega') + (\text{vacuum part}).
\]
This is the \textbf{Bose--Einstein distribution at temperature $T_H$}:
\[
n_H^\mathrm{Planck}(\omega) = \frac{1}{e^{\hbar\omega/k_B T_H} - 1}.
\]
\textbf{Dimensional check:} $\hbar\omega$ has units of energy; $k_B T_H$ has units of energy; ratio is dimensionless; $n_H^\mathrm{Planck}$ is dimensionless (occupation number). $\checkmark$

\subsection{Identification with semiclassical Hawking}

The semiclassical Hawking result (Hawking 1974, 1975) gives the thermal radiation flux at temperature $T_H = \kappa/(2\pi)$ via matter-field-mode tracing through the BH spacetime.

ED's substrate-level KMS derivation yields \textbf{the same leading-order spectrum} via DCGT coarse-graining (Paper\_073 \S6.4): substrate-level V5 KMS coarse-grains to standard matter-field KMS in the continuum, with the same temperature.

\textbf{At leading order} ($\omega \ll \omega_c = c/\ell_P$), ED reproduces standard semiclassical Hawking. The substrate-level mechanism supplies a structural derivation of the temperature via V5 imaginary-time periodicity, but does not change the leading-order spectrum.

\subsection{Why this matters}

Standard semiclassical Hawking requires extrapolation of matter-field modes to $\omega \to \infty$ --- the trans-Planckian problem (Paper\_039 \S5). ED's substrate-level KMS derivation:

\begin{itemize}[noitemsep]
\item Does not require trans-Planckian extrapolation (V5 has finite memory at $\tau_{V5}^\mathrm{BH} \sim \ell_P/c$).
\item Produces structurally-truncated spectrum at $\omega \to \omega_c$.
\item Predicts specific deviations from purely thermal behavior at higher orders (\S6).
\end{itemize}

These three features are the substrate-level structural improvements over standard semiclassical Hawking. The leading-order spectrum coincides; the higher-order content distinguishes.

\section{Non-Thermal Corrections from V5 Cutoff}

\subsection{The V5 substrate cutoff}

By Paper\_090 \S5.2, V5's near-horizon memory time is \textbf{chosen to match the substrate scale}: $\tau_{V5}^\mathrm{BH} = \ell_P/c$ (conditional on $\ell_\mathrm{ED} = \ell_P$, Paper\_027 \S5.3). The corresponding cutoff frequency is:
\[
\omega_c = 1/\tau_{V5}^\mathrm{BH} = c/\ell_P.
\]
V5 envelope $F_{V5}$ decays exponentially for $\omega > \omega_c$; substrate-level V5 correlators cannot support modes above $\omega_c$.

\subsection{Non-thermal correction form --- given P-V5-Even}

The substrate-level Hawking spectrum, derived from V5 KMS with finite cutoff, has the form:
\[
n_H^\mathrm{ED}(\omega) = n_H^\mathrm{Planck}(\omega) \cdot f_{V5}(\omega/\omega_c)
\]
where $f_{V5}(x)$ is the \textbf{V5 cutoff-correction factor} inherited from V5 envelope shape.

\textbf{The leading-order form of $f_{V5}$ depends on V5 envelope parity:}

\textbf{Under P-V5-Even (\S2 postulate):} the V5 envelope is parity-symmetric ($F_{V5}(\omega) = F_{V5}(-\omega)$ around the substrate-saturated regime). The Taylor expansion in $\omega/\omega_c$ contains only \textbf{even powers}, giving:
\[
f_{V5}(x) = 1 - c_1\,x^2 + \mathcal{O}(x^4)
\]
with $c_1 > 0$ (sign-definite, \S6.3). Example: Gaussian $F_{V5}(\omega) \sim e^{-\omega^2/\omega_c^2}$ gives $c_1 = 1/2$.

\textbf{Without P-V5-Even (parity-asymmetric envelopes, e.g., pure exponential $F_{V5}(\omega) \sim e^{-\omega/\omega_c}$):} the Taylor expansion contains a non-zero \textbf{linear term}, giving:
\[
f_{V5}(x) = 1 - c_0\,x + \mathcal{O}(x^2)
\]
with $c_0 \approx 1$ for the exponential case. This is a \textbf{different empirical signature} (linear leading correction rather than quadratic).

\textbf{This paper commits to P-V5-Even} as a substrate-level postulate (\S2). The leading correction is therefore $(\omega/\omega_c)^2$, not $(\omega/\omega_c)$. The choice of even envelope is motivated by Gaussian-like smoothness of the V5 substrate-saturated regime, but it is a \textbf{postulational commitment}, not a derivation --- alternative substrate ontologies with non-even envelopes would predict a different (linear) leading correction.

\textbf{Headline result (under P-V5-Even):}
\[
\boxed{n_H^\mathrm{ED}(\omega) = n_H^\mathrm{Planck}(\omega)\left[1 - c_1\,(\omega/\omega_c)^2 + \mathcal{O}((\omega/\omega_c)^4)\right]}
\]
with $c_1$ order-unity, value-layer (inherited from V5 envelope shape).

\textbf{Dimensional check:} $\omega$ has $[1/\mathrm{time}]$; $\omega_c$ has $[1/\mathrm{time}]$; $\omega/\omega_c$ dimensionless; $n_H$ dimensionless. $\checkmark$

\textbf{Empirical falsification handle.} The choice between $(\omega/\omega_c)^2$ (P-V5-Even) and $(\omega/\omega_c)$ (non-even envelope) is empirically distinguishable in analog-Hawking experiments: linear vs.\ quadratic scaling of the deviation from purely thermal is a sharp test. See \S8 for falsifier sentences.

\subsection{Sign-definite, mass-dependent, information-bearing}

Three structural features distinguish ED's correction from ``tiny correction smoothed away'':

\textbf{Sign-definite (negative).} The leading correction term is $-c_1\,(\omega/\omega_c)^2$ with $c_1 > 0$. The correction \textbf{suppresses emission relative to thermal} at all $\omega < \omega_c$. This is set by V5 envelope decay (positive decay rate at large $\omega$).

\textbf{Mass-dependent.} The dominant Hawking-emission frequency scales as $\omega_H \sim \kappa \sim 1/M$ (for Schwarzschild). The ratio $\omega_H/\omega_c$ scales as $\omega_H/\omega_c \sim M_P/M$. For astrophysical BHs ($M \gg M_P$), $(\omega_H/\omega_c)^2 \sim (M_P/M)^2$ is extremely small --- corrections invisible. For evaporating BHs approaching $M_\star = M_P$ (Paper\_039 \S6), $(\omega_H/\omega_c)^2 \sim 1$ --- corrections $\mathcal{O}(1)$, spectrum qualitatively non-thermal.

\textbf{Information-bearing (soft claim, pending Paper\_050).} Pure thermality carries no information about BH formation state. The substrate framework \textbf{leaves room for} the $(\omega/\omega_c)^2$ deviations to carry V5-cross-chain-correlation content with substrate-level BH-formation history. This is the substrate-level \emph{channel} through which BH-formation information could be carried out by the radiation (Paper\_039 \S7.2 + Paper\_050 in queue).

\textbf{Honesty caveat:} the explicit map from substrate-level BH-formation state to specific spectrum features is \textbf{not derived in this paper}. ``Information-bearing'' is a substrate-level \emph{room-for}, not a derivation. The explicit encoding map awaits Paper\_050 (Page-curve substrate mechanism). Without that map, ``information-bearing'' is a structural-availability claim, not an operational identification. Substrate-level unitarity is established structurally in Paper\_039 \S7; this paper supplies the substrate-level \emph{vehicle} (non-thermal corrections) through which information \emph{can be} carried, not the explicit encoding.

\subsection{Empirical content}

\begin{itemize}[noitemsep]
\item \textbf{Astrophysical BHs:} corrections at level $(\hbar\omega/M_P c^2)^2 \sim (T_H/T_P)^2$ --- extremely small but structurally non-zero. Direct astrophysical detection beyond current sensitivity.
\item \textbf{Analog Hawking experiments} (BEC, fluid, optical analogs): analog system's natural high-frequency cutoff much closer to dominant analog-Hawking frequency. Non-thermal corrections should be observable as deviation from purely thermal prediction. \textbf{Near-term empirical test.}
\item \textbf{Late-stage BH evaporation:} at $M \to M_\star = M_P$, corrections become $\mathcal{O}(1)$; spectrum qualitatively non-thermal. Direct observation in PBH-evaporation cosmological signatures (if any) or hypothetical near-Planck-mass BH evaporation.
\end{itemize}

\subsection{Coefficient $c_1$ value-layer status}

$c_1$ depends on V5 envelope shape (Paper\_090 \S3.3):

\begin{center}
\begin{tabular}{ll}
\toprule
\textbf{Envelope shape} & $c_1$ \\
\midrule
Exponential $e^{-\omega/\omega_c}$ & $\approx 1$ \\
Gaussian $e^{-\omega^2/\omega_c^2}$ & $1/2$ \\
Step function $\Theta(\omega_c - \omega)$ & divergent at $\omega = \omega_c$ (not physical) \\
\bottomrule
\end{tabular}
\end{center}

The specific value of $c_1$ is \textbf{value-layer empirical}: it requires independent determination of V5 envelope shape (e.g., via cross-domain V5 applications in soft-matter Maxwell relaxation, where empirical relaxation kernels constrain V5 envelope independently).

The substrate framework's structural prediction is the \textbf{form} ($(\omega/\omega_c)^2$ scaling, negative sign at leading order, mass-dependence); the \textbf{coefficient} is downstream of V5 envelope determination.

\section{The Wedge --- Empirical Distinction from Standard Semiclassical Hawking}

The substrate-level Hawking spectrum is \textbf{distinguishable from purely thermal semiclassical Hawking} via:

\begin{enumerate}[noitemsep]
\item \textbf{Sign-definite negative corrections at $(\omega/\omega_c)^2$.} Standard semiclassical Hawking is exactly thermal; ED predicts specific deviations.
\item \textbf{Mass-dependent scaling of the correction magnitude.} Standard semiclassical Hawking has no analogous correction; ED predicts that the deviation grows as $M$ decreases.
\item \textbf{Information-bearing non-thermal content.} Standard semiclassical Hawking is information-less; ED's non-thermal corrections carry substrate-level information about BH formation state.
\end{enumerate}

\textbf{Empirical strategy:}

\begin{itemize}[noitemsep]
\item Analog Hawking experiments (closest near-term test): measure the analog spectrum's deviation from purely thermal at the analog system's $(\omega/\omega_c^\mathrm{analog})^2$ scale.
\item Late-stage BH evaporation (long-term test): direct astrophysical observation if PBH evaporation produces detectable signatures.
\end{itemize}

This is the wedge content of Paper\_047: a \textbf{specific, falsifiable, sign-definite, mass-dependent non-thermal correction} that no standard semiclassical or modified-gravity account predicts.

\section{Falsification Criteria}

\subsection{Analog Hawking with exact thermal spectrum at $(\omega/\omega_c)^2$}

\textbf{Falsifier sentence:} \emph{Analog Hawking experiments (BEC, fluid, optical) detecting Hawking-like radiation with no detectable deviation from purely thermal at the analog system's $(\omega/\omega_c^\mathrm{analog})^2$ scale, after accounting for analog-system noise, would falsify \S6's headline non-thermal-correction prediction.}

\subsection{Wrong sign or scaling of non-thermal corrections (sharp P-V5-Even test)}

\textbf{Falsifier sentence (sign):} \emph{Empirical observation of non-thermal Hawking-spectrum corrections with positive sign (enhancement rather than suppression) would falsify V5 envelope decay shape.}

\textbf{Falsifier sentence (parity / scaling):} \emph{Empirical observation of a linear $(\omega/\omega_c)^1$ leading correction (rather than quadratic $(\omega/\omega_c)^2$) would falsify the P-V5-Even postulate (\S2) and force adoption of a parity-asymmetric V5 envelope. This is a sharp empirical test of P-V5-Even directly: the leading-correction power-law is the falsification handle.}

\textbf{Falsifier sentence (higher-order):} \emph{Empirical scaling other than $(\omega/\omega_c)^2$ or $(\omega/\omega_c)^1$ (e.g., $(\omega/\omega_c)^{1.5}$, or no power-law scaling at all) would falsify the V5 envelope-Taylor-expansion mechanism entirely, requiring a different substrate-level mechanism for the non-thermal corrections.}

\subsection{No imaginary-time periodicity at substrate level}

\textbf{Falsifier sentence:} \emph{Substrate-level analysis showing V5 cross-chain correlators do not acquire imaginary-time periodicity in the near-horizon Euclidean substrate-graph structure would falsify \S3.3 and the KMS / Hawking-temperature derivation.}

\subsection{Hawking temperature differs from $\kappa/(2\pi)$}

\textbf{Falsifier sentence:} \emph{Astrophysical observation of BH thermal radiation at temperature substantially differing from $\kappa/(2\pi)$ --- beyond the predicted $(\omega/\omega_c)^2$ corrections --- would falsify the substrate-level KMS mechanism.}

\subsection{V5 cutoff frequency differs from $c/\ell_P$}

\textbf{Falsifier sentence:} \emph{Empirical demonstration that the trans-Planckian cutoff for the Hawking spectrum is at a frequency $\omega_c \neq c/\ell_P$ --- i.e., inconsistent with $\tau_{V5}^\mathrm{BH} = \ell_P/c$ matching --- would falsify the substrate-scale matching of Paper\_090 \S5.2 (which Paper\_047 inherits).}

\subsection{Late-stage evaporation purely thermal at $M \to M_P$}

\textbf{Falsifier sentence:} \emph{Empirical evidence that late-stage BH evaporation near $M = M_P$ produces purely thermal spectrum (no $\mathcal{O}(1)$ non-thermal deviations as predicted by \S6.3 mass-dependence) would falsify the substrate-level $(\omega/\omega_c)^2$ correction structure.}

\section{Position Statement}

The \textbf{Hawking spectrum at $T_H = \kappa/(2\pi)$} is a downstream structural identification in the ED Generative Primitives System. Given the postulated primitives + V1/V5 + Paper\_039's substrate decoupling-surface mechanism + Paper\_073 DCGT, the spectrum follows from:

\begin{itemize}[noitemsep]
\item V5 imaginary-time periodicity $\beta_H = 2\pi/\kappa$ at the Euclidean near-horizon substrate-graph structure (\S3).
\item KMS condition on V5 cross-chain correlators (\S4).
\item Thermal Planck distribution at $T_H = 1/\beta_H = \kappa/(2\pi)$ (\S5).
\end{itemize}

\textbf{Headline non-thermal correction} (\S6, delivering Paper\_039 \S7.2's promise):
\[
n_H^\mathrm{ED}(\omega) = n_H^\mathrm{Planck}(\omega)[1 - c_1\,(\omega/\omega_c)^2 + \mathcal{O}((\omega/\omega_c)^4)]
\]
with:

\begin{itemize}[noitemsep]
\item $\omega_c = c/\ell_P$ (V5 substrate cutoff, chosen to match substrate scale).
\item $c_1$ order-unity, value-layer (V5 envelope shape).
\item Sign-definite negative, mass-dependent, information-bearing.
\end{itemize}

What this paper claims:

\begin{itemize}[noitemsep]
\item Given the postulated primitives + V5 + Paper\_039 + Paper\_073, $T_H = \kappa/(2\pi)$ is the unique substrate-level identification of the Hawking temperature.
\item The leading-order spectrum coincides with standard semiclassical Hawking.
\item The $(\omega/\omega_c)^2$ non-thermal correction is a specific, falsifiable, sign-definite prediction that distinguishes ED from standard semiclassical accounts.
\item Analog Hawking experiments are the near-term empirical wedge.
\end{itemize}

What this paper does \emph{not} claim:

\begin{itemize}[noitemsep]
\item The substrate primitives are not derived (Paper\_087).
\item $\kappa$ is not derived --- inherited from BH geometry.
\item $\omega_c$ is not derived --- chosen to match substrate scale ($\ell_\mathrm{ED} = \ell_P$).
\item $c_1$ is not exactly derivable --- value-layer V5 envelope shape.
\item Unitarity proof not delivered in this paper (Paper\_039 + Paper\_052 in queue).
\item BH mass/geometry not derived.
\end{itemize}

The empirical case is cross-domain reach + the analog-Hawking-experiment near-term test.

\textbf{Series context.} Paper\_047 of the BH/Hawking arc. The arc continues:

\begin{itemize}[noitemsep]
\item \textbf{Paper\_050 (in queue):} Page-curve substrate mechanism.
\item \textbf{Paper\_052 (in queue):} BH info-paradox publication-grade synthesis.
\end{itemize}

\section*{Appendix: Cross-References and Glossary}

\subsection*{Cross-references}

\begin{itemize}[noitemsep]
\item \textbf{Position paper:} Paper\_087.
\item \textbf{Paper\_089 (V1):} kernel inheritance.
\item \textbf{Paper\_090 (V5):} cross-chain kernel; cutoff $\omega_c$ matching.
\item \textbf{Paper\_039 (Horizon Decoupling):} substrate decoupling surface; $\kappa$ inheritance.
\item \textbf{Paper\_073 (DCGT):} continuum-limit bridge.
\end{itemize}

\subsection*{Glossary}

\begin{itemize}[noitemsep]
\item \textbf{Hawking temperature $T_H = \kappa/(2\pi)$.} Substrate-level identification via V5 KMS.
\item \textbf{Surface gravity $\kappa$.} Inherited from BH geometry; not derived in this paper.
\item \textbf{Imaginary-time periodicity $\beta_H$.} $\beta_H = 2\pi/\kappa$ from Euclidean conical-defect smoothness.
\item \textbf{KMS condition.} Kubo--Martin--Schwinger thermal-state condition on V5 correlators.
\item \textbf{V5 substrate cutoff $\omega_c = c/\ell_P$.} Chosen to match substrate scale.
\item \textbf{Non-thermal correction.} $(\omega/\omega_c)^2$ leading-order deviation from purely thermal; sign-definite, mass-dependent, information-bearing.
\item \textbf{Coefficient $c_1$.} Order-unity, value-layer V5 envelope shape.
\item \textbf{Conical-defect smoothness.} Substrate-graph regularity condition at Euclidean horizon.
\end{itemize}

\end{document}
